% Generated draft: requires mathematical and source review.
\documentclass[11pt]{article}
\usepackage[margin=25mm]{geometry}
\usepackage{fontspec}
\setmainfont{Latin Modern Roman}
\usepackage{amsmath,amssymb,mathtools,unicode-math}
\setmathfont{Latin Modern Math}
\usepackage{enumitem,needspace,xurl,hyperref}
\setlength{\parindent}{0pt}
\setlength{\parskip}{5pt}
\setlength{\emergencystretch}{3em}
\newcommand{\N}{\mathbb{N}}
\newcommand{\Z}{\mathbb{Z}}
\newcommand{\Q}{\mathbb{Q}}
\newcommand{\R}{\mathbb{R}}
\newcommand{\C}{\mathbb{C}}
\newcommand{\F}{\mathbb{F}}
\newcommand{\A}{\mathbb{A}}
\newcommand{\PP}{\mathbb{P}}
\newcommand{\E}{\mathbb{E}}
\newcommand{\eps}{\varepsilon}
\newcommand{\dd}{\,\mathrm{d}}
\begin{document}
\section*{UnsolvedMath ID: 30004527}
\textbf{Source problem code:} catalogue entry
\subsection{Definitions and mathematical statement}
A conventional interpretation of the three-dimensional question can be formulated as follows. Fix an algebraically closed field \(k\). This is an explicit provisional convention: the supplied record specifies neither the ground field nor its characteristic.

Let \(Z\subset\mathbb{P}^{3}_{k}\) be a finite reduced set of points. For positive integers \(a,b\), say that \(Z\) has the \((a,b)\)-geproci property if its projection from a general center \(p\in\mathbb{P}^{3}_{k}\) to a plane \(H\cong\mathbb{P}^{2}_{k}\) is injective on \(Z\) and its reduced image is a scheme-theoretic complete intersection
\[
\pi_p(Z)=V_H(F_p,G_p),\qquad \deg F_p=a,\quad \deg G_p=b,
\]
where \(F_p,G_p\) have no common curve component. Here ``general'' means that the condition holds for every center in some nonempty Zariski open subset of \(\mathbb{P}^{3}_{k}\setminus(H\cup Z)\). In particular, \(|Z|=ab\). Say that \(Z\) is geproci if it has this property for some \(a,b\).

Under a positive-integer convention, an \((r,s)\)-grid is a set
\[
Z=\{L_i\cap M_j:1\leq i\leq r,\ 1\leq j\leq s\},
\]
where the \(L_i\) are pairwise skew lines, the \(M_j\) are pairwise skew lines, every \(L_i\) meets every \(M_j\), and the \(rs\) intersection points are distinct. Call \(Z\) non-grid if it admits no such representation. The supplied record does not define grids or specify whether cases with \(r=1\) or \(s=1\), or geproci degrees equal to \(1\), are admitted; this convention therefore requires confirmation.

With these provisional conventions, the first question is whether
\[
\forall N\in\N\ \exists a,b\in\Z_{\geq1}\ \exists Z\subset\mathbb{P}^{3}_{k}:
\quad Z\text{ is finite and reduced},\quad
|Z|=ab>N,\quad
Z\text{ is non-grid and }(a,b)\text{-geproci}.
\]
The field \(k\) is a parameter of this formulation; the source does not say whether existence is sought over one particular field or uniformly over a class of fields. The word ``family'' is interpreted here only as unbounded cardinality, without imposing an algebraic parameter space.

The second question asks for analogous examples in higher-dimensional projective spaces. The supplied record does not specify the projection target, the complete-intersection condition, a higher-dimensional definition of grid, or whether examples are required in some or every dimension greater than three. Consequently, no unique quantified higher-dimensional statement can be extracted.

The supplied dated curation labels the problem partially solved and reports non-grid geproci examples, including infinite families and a \(24\)-point example, while describing the requested general series and higher-dimensional analogue as unresolved. These are supplied assessments, not independently verified conclusions; the relationship between the reported infinite families and the displayed unbounded-cardinality question is not explained.
\subsection{Short English statement}
The question asks whether non-grid finite point sets in projective three-space can have arbitrarily large cardinality while their general plane projections are complete intersections, and whether a comparable phenomenon occurs in higher dimensions. The ground field, degeneracy conventions, and higher-dimensional meaning require clarification. The supplied partially solved label is dated curation rather than an independently verified status.
\subsection{Sources}
\begin{itemize}
\item Supplied local UnsolvedMath registry record: permanent ID \(30004527\), title \emph{Unbounded Non-Grid Geproci Point Configurations}, exact source problem code \texttt{catalogue entry}, in the Oberwolfach Reports Open Problems collection. The record gives proposed year \(2020\) and names no proposer.
\item Supplied report citation: \emph{Lefschetz Properties in Algebra, Geometry and Combinatorics} (2021), Oberwolfach Reports, pp.\ 1539--1600, DOI \texttt{10.4171/owr/2020/31}. \url{https://doi.org/10.4171/owr/2020/31}. The background describes these reports as documenting mathematical workshops and their open-problem discussions.
\item Supplied literature reference labeled ``Chiantini--Migliore'': \url{https://doi.org/10.1090/tran/8290}. No title, theorem number, or specific supporting passage is supplied.
\item Literature assessment checked at \texttt{2026-08-22}: supplied dated curation, not independently verified. It labels the problem partially solved, reports infinite families of non-grid geproci examples and a \(24\)-point example in projective three-space, and describes the requested general series and higher-dimensional analogue as unresolved. The background expressly requests independent verification of time-sensitive status and cited claims.
\item The registry labels the statement \texttt{corrected\_verified} and says its mathematical content was preserved while notation, typography, omitted definitions, or duplicate clauses were repaired. This is supplied registry metadata, not an independent mathematical verification of this draft.
\end{itemize}
\end{document}
